\documentclass[10pt,aspectratio=169]{beamer}

% --- Theme ---
% Requires the `metropolis` Beamer theme (texlive-fonts-extra on Linux,
% available in MacTeX). Fallback: comment the next two lines and uncomment
% \usetheme{Boadilla} below.
\usetheme{metropolis}
\metroset{progressbar=none, numbering=fraction, block=fill}
% \usetheme{Boadilla}

% --- Packages ---
\usepackage[utf8]{inputenc}
\usepackage{graphicx}
\usepackage{booktabs}
\usepackage{tikz}
\usetikzlibrary{shapes.geometric, arrows.meta, positioning, calc}
\usepackage{xcolor}
\usepackage{amsmath}
\usepackage{amssymb}
\usepackage{bookmark}
\usepackage{tcolorbox}

% --- Image search paths ---
\graphicspath{
  {../notebooks/images/}
  {../notebooks/images/2025-01-15/}
}

% --- Colors for quadrant labels ---
\definecolor{good}{RGB}{44,160,44}
\definecolor{bad}{RGB}{214,39,40}
\definecolor{warn}{RGB}{255,140,0}

% --- Title metadata ---
\title{Theme Detection - updates}
\date{May 2026}

\begin{document}

% =====================================================================
\begin{frame}[shrink=5]
  \titlepage
  \vspace{-1.0cm}
  \centering\footnotesize
  Goal: surface accelerating investment themes by combining news and ETF signals.\\
  Today we want your input on the design trade-offs.
\end{frame}


% =====================================================================
% Topic detection (3 slides)
% =====================================================================
\begin{frame}[t]{Topic detection: classical methods (1/2)}
  \footnotesize
  \begin{columns}[T,onlytextwidth]
    \column{0.52\textwidth}
      \textbf{Bag of Words.}\;
      Documents $\rightarrow$ vocabulary $\rightarrow$ document--term matrix $\mathbf{X}$.
      \vspace{0.15cm}

      \textbf{Example docs:}
      {\ttfamily\scriptsize
        cats chase mice \textbar\ dogs chase cats \textbar\ stocks and markets rise}

      \vspace{0.2cm}
      \textbf{Hard assignment} (K-means, hierarchical clustering): one label per doc\\
      (D1, D2 $\rightarrow$ Animals; D3 $\rightarrow$ Finance).
      \emph{Limitation:} real text is multi-topic.

      \vspace{0.2cm}
      \textbf{Matrix factorization:}\;
      $\mathbf{X} \approx \mathbf{W}\mathbf{H}$
      ($\mathbf{W}$ doc--topic, $\mathbf{H}$ topic--word) $\Rightarrow$ soft mixtures, latent structure.
      Methods: LSA (SVD), NMF.

    \column{0.46\textwidth}
      \centering
      \textbf{Document--term matrix}
      \vspace{0.1cm}

      {\scriptsize
      \begin{tabular}{@{}lcccc@{}}
        \toprule
        Doc & cats & chase & mice & stocks \\
        \midrule
        D1 & 1 & 1 & 1 & 0 \\
        D2 & 1 & 1 & 0 & 0 \\
        D3 & 0 & 0 & 0 & 1 \\
        \bottomrule
      \end{tabular}}

      \vspace{0.35cm}
      \textbf{pLSA / LDA} (probabilistic mixtures):
      \[
        P(w \mid d) = \textstyle\sum_t P(w \mid t)\, P(t \mid d)
      \]
      LDA adds Bayesian priors: each document has a distribution over topics; each topic over words.
      Became the classical default (interpretable, scalable).
  \end{columns}
\end{frame}

\begin{frame}[t]{Topic detection: embeddings and BERTopic (2/2)}
  \footnotesize
  \textbf{Shift:} word counts $\rightarrow$ semantic vectors (Word2Vec, BERT, Sentence-BERT).
  Each document $d_i \in \mathbb{R}^{768}$; similar meaning $\Rightarrow$ nearby in space.

  \vspace{0.25cm}
  \textbf{BERTopic pipeline:}
  documents $\rightarrow$ embeddings $\rightarrow$ UMAP $\rightarrow$ HDBSCAN $\rightarrow$ c-TF-IDF keywords.
  Combines transformer embeddings, nonlinear dimensionality reduction, density clustering, and class-based TF-IDF for interpretable labels.

  \vspace{0.25cm}
  \textbf{No fixed $K$:} topics are \emph{dense regions} in embedding space---HDBSCAN finds stable clusters of semantically similar documents rather than prescribing the number of topics upfront.


  \begin{columns}[T,onlytextwidth]
    \column{0.48\textwidth}
      \begin{block}{Classical}
        \footnotesize
        Topics are \textbf{word frequency distributions} inferred from $\mathbf{X}$:\\
        hard labels, factorized latent factors, or generative models (pLSA, LDA).
      \end{block}

    \column{0.48\textwidth}
      \begin{block}{Modern (BERTopic)}
        \footnotesize
        Topics are \textbf{regions in semantic embedding space}:\\
        geometry of meaning drives clustering; keywords summarize each region post hoc.
      \end{block}
  \end{columns}

\end{frame}

% =====================================================================
\begin{frame}[t]{Open questions: what is a theme?}
  \small
  Before we fix the pipeline, we need shared answers (or explicit trade-offs) on:

  \begin{enumerate}
    \item What is a \textbf{theme}?
    \item Is a theme the same as a \textbf{topic} (in the NLP sense)?
    \item Is a theme a \textbf{cluster of news}---or an \textbf{abstraction} on top of a cluster?
    \item Does it have a clear \textbf{geographical} location or scope?
    \item Does it have a clear \textbf{temporal} placement or scope?
    \item What \textbf{level of granularity} should it have?
    \item For each theme we surface: can it be \textbf{monetized} (mapped to investable exposure)?
  \end{enumerate}

\end{frame}


\end{document}
